\documentclass[11pt]{article}
\usepackage[margin=1in]{geometry}
\usepackage{amsmath,amssymb}
\usepackage{booktabs}
\usepackage[hidelinks]{hyperref}
\usepackage{kotex}
% \usepackage{microtype} % disabled: bitmap Hangul fonts break font expansion

\title{How to Train a Jev-Like Model\\
\large Treating Typed Evaluation as One-Token Classification}
\author{Jioh Jung\\ \texttt{ziozzang@skku.edu}}
\date{September 21, 2026}

\begin{document}
\maketitle

\begin{abstract}
A Jev-like model does not need to learn to write an answer.  It needs to
place probability mass on one member of a declared finite answer set and to
do so at the first decode position.  This observation turns model training
from long-horizon language generation into classification with a language
model head.  When reinforcement learning (RL) is appropriate, the episode
has one action and one terminal reward: emit one label token.  The resulting
problem is a contextual bandit, avoiding credit assignment over reasoning
traces, verbosity rewards, answer extraction, and most format failures.
The complete class distribution remains available from the label logits,
so a sampled token is used only for training and the serving path need not
sample at all.  We describe the objective, label-token design, reward and
data construction, calibration, and deployment checks.  We also state the
important negative result: tuning is optional.  A general instruction model
that already ranks the choices correctly can be used unchanged by reading
its next-token log-probabilities.  Fine-tuning is primarily justified when
the deployment requires a more stable first-token interface, a domain-specific
decision boundary, or improved calibration---not merely to make the model
look specialized.
\end{abstract}

\section{The training target is a classifier}

The companion hearim paper~\cite{hearim} shows how an ordinary
autoregressive model can implement the Jev \texttt{/v1/systemone} contract:
compile a typed question into a labeled choice, read the candidate labels'
next-token log-probabilities, and conditionally normalize them.  If
$h_\theta(x)$ is the model state after prompt $x$ and $z_\theta(\ell\mid x)$
is the logit for label token $\ell$, then the answer distribution is
\begin{equation}
q_\theta(k\mid x,A) =
\frac{\exp z_\theta(\ell_k\mid x)}
     {\sum_{j\in A}\exp z_\theta(\ell_j\mid x)},
\label{eq:conditional}
\end{equation}
where $A$ is the set of answers declared by the request.  Choice, ordinal
score, and boolean (\emph{noul}) questions are all reductions of this one
categorical distribution.

Equation~\ref{eq:conditional} is also the correct training abstraction.
The model is a classifier whose class head happens to be selected rows of
its language-model vocabulary matrix.  Training it to produce explanations,
JSON, or a chain of thought solves a larger problem than the API asks for
and introduces irrelevant failure modes.  The desired completion is exactly
one token; EOS may follow, but it is outside the decision itself.

This design has a practical precedent.  Kakao's Kanana Safeguard 8B is a
causal language model trained to classify a user utterance or assistant
response and emit a single token such as \texttt{<SAFE>} or
\texttt{<UNSAFE-S4>}; its published inference example sets
\texttt{max\_new\_tokens=1}, and its evaluation uses the first output
token~\cite{kanana}.  A Jev-like model generalizes the same interface from a
fixed safety taxonomy to a request-defined taxonomy.

KT's Llama SafetyGuard Content Binary model is an even closer example of the
probability interface: it declares itself a text classifier, reserves the
single tokens \texttt{<SAFE>} and \texttt{<UNSAFE>}, serves one token, and
applies a configurable threshold to their two logits~\cite{contentbinary}.
Its streaming deployment repeatedly classifies the growing response prefix;
this changes the input schedule, not the one-token decision.  Both guard
models are fixed-taxonomy classifiers, whereas a Jev-like model receives the
candidate meanings in each request.  The distinction affects data diversity
and label permutation, but not the central architecture: a causal LM can act
as a classifier by making its first-token vocabulary rows the class head.

\section{Labels and prompt contract}

Let the serving prompt end immediately before the label position:
\begin{verbatim}
State: ...
Question: Which handling policy applies?
<answer-a> = refund
<answer-b> = replace
<answer-c> = escalate
Answer:
\end{verbatim}
The literal spelling is not important.  The following invariants are.

\begin{enumerate}
\item Every label must be exactly one token in the answer context.  Verify
tokenization after the delimiter; do not infer it from how the label looks.
\item The entire candidate alphabet must be available at training and
serving time.  Labels are assigned independently of answer meaning and their
order is permuted during training to prevent a semantic prior from attaching
to a position.
\item The production chat template, system text, delimiter, and answer marker
must match training.  The loss is applied only at the first answer position.
\item Inference reads all candidate logits.  It does not estimate a
distribution by repeatedly sampling the one-token answer.
\end{enumerate}

Existing stable tokens such as digits or letters are usually sufficient and
require no tokenizer or embedding change.  Dedicated tokens such as
\texttt{<answer-a>} are optional.  They are useful when no sufficiently large
alphabet remains single-token at the relevant boundary, when label strings
collide with normal-language priors, or when a hard protocol boundary is
operationally valuable.  Otherwise they add tokenizer migration, embedding
initialization, checkpoint compatibility, and serving configuration work
without changing the statistical problem.  If added, the tokenizer and
model vocabulary must be resized together, every new token must receive a
trainable embedding/output row, and the complete deployment stack must be
verified against that exact tokenizer revision.

\section{Choose the least training that solves the problem}

There are three increasingly invasive regimes.

\paragraph{No tuning.}
First measure an ordinary instruction model using the exact serving prompt.
If class accuracy, negative log-likelihood, calibration, candidate mass, and
label-order stability meet the product target, use it unchanged.  The model
already computed the needed distribution; changing its weights creates no
intrinsic benefit.

\paragraph{Supervised one-token fine-tuning.}
When a gold class $y$ is available for every example, masked next-token
cross-entropy is the direct objective
\begin{equation}
\mathcal{L}_{\mathrm{SFT}}(\theta)
= -\log q_\theta(y\mid x,A).
\label{eq:sft}
\end{equation}
It is normally simpler and more sample-efficient than on-policy RL for the
same fixed labels.  Training examples should contain the production prompt,
one answer token, and no rationale; all prompt-token losses are masked.  SFT
is the right first tool when the requested behavior can be written down as a
labeled classification corpus.

\paragraph{One-step RL.}
RL is attractive when correctness is supplied by an executable verifier,
policy rule, downstream outcome, preference comparison, or other reward that
is easier to evaluate than to label densely.  It is also a natural final
alignment stage after SFT.  Crucially, a Jev-like episode does not need a
reasoning trajectory:
\[
x \longrightarrow a\in A \longrightarrow r(x,a).
\]
There is one state, one sampled label token, and one terminal reward.  This
is a contextual bandit.  Policy-gradient training therefore has no temporal
credit-assignment problem:
\begin{equation}
\nabla_\theta J =
\mathbb{E}_{x,a\sim q_\theta}
\left[(r(x,a)-b(x))\nabla_\theta\log q_\theta(a\mid x,A)\right],
\label{eq:pg}
\end{equation}
where $b(x)$ is a learned value, a group-relative baseline, or the mean
reward of several sampled labels.  PPO~\cite{ppo} or a group-relative update
can implement this objective, but the important optimization is not the
algorithm's name: it is deleting every action after the first token.

For a correct class $y$, a minimal reward is
\begin{equation}
r(x,a)=\mathbf{1}[a=y]-\lambda_{\mathrm{invalid}}
\mathbf{1}[a\notin A].
\end{equation}
Cost-sensitive classification can substitute a reward matrix $R(y,a)$;
ordinal questions can reward near misses smoothly; a noul reward can be a
proper score from a delayed binary outcome.  When probability quality matters,
a proper scoring rule should also be optimized or used for model selection.
Rewarding only argmax correctness does not by itself calibrate the remaining
logits.

The update should include a modest KL penalty to a reference model when
preserving general behavior matters.  Candidate-masked sampling makes every
action syntactically valid and improves exploration, but it changes the
training policy's probability space.  Deployments that score unmasked logits
must additionally penalize probability mass outside $A$ or mix in examples
trained on the full vocabulary.  Otherwise a model may classify correctly
only when a decoder grammar rescues it.

\section{Efficient recipe}

An efficient training run follows a short loop.

\begin{enumerate}
\item Freeze the production prompt contract and construct a stable
single-token alphabet large enough for the maximum class count.
\item Build domain examples as $(x,A,y)$ or $(x,A,R)$ and split them by
underlying entity or scenario, not by surface paraphrase.  Randomize label
assignment and answer order independently for every presentation.
\item Establish the untuned baseline.  Record accuracy, macro-F1, NLL, Brier
score, expected calibration error, candidate mass, invalid-token mass, and
permutation variance.
\item If explicit labels exist, perform masked one-position SFT.  Parameter-
efficient adapters can reduce memory and make rollback simple; full tuning is
not required by the interface.
\item If rewards are programmatic, preference-based, or downstream, continue
with one-step RL.  Sample only the label position, compute terminal reward,
and update against a frozen reference with KL control.  Mix difficult and
ambiguous examples rather than spending compute generating rationales.
\item Calibrate on held-out data.  A scalar temperature $T$ fitted to NLL
gives $q_T(k)\propto\exp(z_k/T)$ without changing argmax accuracy.  Maintain
separate calibration profiles for every model, tokenizer, prompt template,
and scoring mode.
\item Reject the new checkpoint unless it improves the target metric while
passing token-boundary, label-permutation, invalid-mass, and out-of-domain
regression tests.
\end{enumerate}

For $K$ classes, naive on-policy exploration becomes sparse as $K$ grows.
If the correct answer is already known, Equation~\ref{eq:sft} uses that
information exactly and should not be replaced by sampling.  If only a
black-box reward is available, multiple label samples or enumerating all
$K$ label rewards can reduce variance; when all rewards can be evaluated,
the expected reward $\sum_k q_\theta(k\mid x,A)R(x,k)$ can be differentiated
directly.  Calling the latter procedure ``RL'' is organizationally harmless,
but computationally it has recovered a full-information classification
update.

\section{What tuning is---and is not---for}

Fine-tuning cannot make a serving API reveal logits it withholds, cannot
repair a tokenizer mismatch outside the trained context, and cannot turn an
argmax-only endpoint into a calibrated probability service.  Those are
interface constraints.  Nor should fine-tuning be used merely to force a
model to narrate less: raw completions, constrained decoding, or direct
logit access may already provide the needed surface.

The defensible reasons to tune are narrower:
\begin{itemize}
\item raise task accuracy on a specific decision boundary;
\item concentrate answer probability at the first position and reduce mass
on prose, reasoning openers, or malformed labels;
\item reduce sensitivity to label permutation and prompt paraphrase;
\item fit domain costs or rewards not represented in generic instruction
tuning; and
\item improve calibration, validated on held-out examples.
\end{itemize}
Thus output stability is the default justification.  Specialization is an
option when measurements demand it, not a prerequisite for Jev-compatible
serving.

\section{Evaluation and release gate}

Report more than first-token accuracy.  Macro-F1 detects neglected classes;
NLL and Brier score test the full distribution; ECE gives a coarse calibration
view; candidate mass reveals whether the model prefers non-label text; and
label-order permutation variance exposes positional shortcuts.  Slice all
metrics by question type, class count, language, domain, ambiguity, and
prompt length.  Evaluate both constrained and unconstrained decoding, while
using unconstrained candidate logits for the published Jev distribution.

The release artifact must bind the weights to a tokenizer digest, chat or
completion template, delimiter, label alphabet, and calibration temperature.
A startup probe should verify that each label is one token in context and
that the server returns every candidate log-probability.  These checks matter
even for a perfectly trained checkpoint: the output protocol is partly a
property of the serving stack.

\section{Conclusion}

The most economical Jev-like model is a classifier presented through an
autoregressive interface.  If training is needed, supervise or reinforce one
decision token, not a generated explanation.  One-step RL is especially
clean when reward is easier to obtain than a gold label: it reduces the
episode to a contextual bandit and makes format compliance part of the same
terminal decision.  Yet no-training remains a first-class result.  General
models often rank short labeled choices well enough already, and hearim can
read that distribution without changing the model.  Dedicated
\texttt{<answer-a>} tokens, SFT, and RL are tools for measured deficiencies---
above all output stability---rather than requirements of the architecture.

\begin{thebibliography}{9}

\bibitem{hearim}
J. Jung.
\newblock \emph{hearim: Turning Generation APIs into Probability Oracles---A
Jev-Compatible Reimplementation of Single-Token Label Scoring}.
\newblock 2026. Companion paper in this repository.

\bibitem{kanana}
Kanana Safeguard Team.
\newblock \emph{Kanana Safeguard}.
\newblock Kakao, 2025.
\newblock \url{https://huggingface.co/kakaocorp/kanana-safeguard-8b}.

\bibitem{contentbinary}
W. Lee, Y. Kim, Y. Park, J. Moon, D. Jeong, and W. Park.
\newblock Guard Vector: Beyond English LLM Guardrails with Task-Vector
Composition and Streaming-Aware Prefix SFT.
\newblock \emph{arXiv:2509.23381}, 2025.
\newblock Model card: \url{https://huggingface.co/K-intelligence/Llama-SafetyGuard-Content-Binary}.

\bibitem{ppo}
J. Schulman, F. Wolski, P. Dhariwal, A. Radford, and O. Klimov.
\newblock Proximal Policy Optimization Algorithms.
\newblock \emph{arXiv:1707.06347}, 2017.

\bibitem{bandit}
T. Lattimore and C. Szepesv\'ari.
\newblock \emph{Bandit Algorithms}.
\newblock Cambridge University Press, 2020.

\end{thebibliography}

\end{document}
